\section{Evaluation}
\label{sec:evaluation}

\subsection{Methodology}

The order of this evaluation is deliberate: every sample was first
analyzed manually and statically, and only afterward was the pipeline
run against it.

Ground truth for each of the three families was produced by manual
reverse engineering, documented as a full analyst report and never
derived from the pipeline's own output. This independence is a
deliberate methodological safeguard, not an incidental detail: both
engines depend on third-party tooling (YARA, FLOSS, CAPE's community
signatures) that can itself mis-map a technique, so ground truth had to
be established without consulting the pipeline's own extracted
artifacts, otherwise validating the pipeline against it would be
circular. Concretely, the same protocol was applied to all three
families, in this order, and only the last step ever touches the
pipeline's own output:

\begin{enumerate}
  \item \textbf{Independent static disassembly.} Every sample was
    loaded into Ghidra and decompiled function by function, each
    function of interest renamed to reflect its established purpose,
    without consulting the pipeline's own import or string extraction
    output.
  \item \textbf{Independent dynamic observation.} Sandbox detonation
    was watched directly (API call sequence, dropped files, network
    activity), noting behavior from the raw trace rather than through
    CAPE's own community signature tags, which \S\ref{sec:dynamic}
    documents as themselves needing curation.
  \item \textbf{Independent configuration and key recovery.} Where a
    sample encrypted or obfuscated its configuration, the key or
    algorithm was recovered manually, entirely separate from the
    pipeline's own automated brute-force and XOR-recovery modules.
  \item \textbf{Only then, the diff.} The pipeline was run and its
    \code{attck\_mapping.json} compared against the finished,
    independently-produced report's own ATT\&CK Techniques table (the
    false-positive audit, \S\ref{sec:fp-audit}). By construction, every
    discrepancy it finds is a disagreement between two
    independently-produced artifacts, not the pipeline checked against
    its own assumptions.
\end{enumerate}

Table~\ref{tab:validation} gives one concrete example per family of
evidence this protocol recovered that the automated pipeline never
touches, which is what makes each report an independent check rather
than a restatement of what the tooling already assumes.

\begin{table}[h]
\centering
\small
\begin{tabular}{@{}lp{9.5cm}@{}}
\toprule
\textbf{Sample} & \textbf{Independently recovered, pipeline never touches it} \\
\midrule
Akira & 17-row ATT\&CK table built from function-level Ghidra decompilation of the sample's threading, logging, and encryption internals. \\
\rowline
RoningLoader & All four chain components (installer, \code{D3D11InstallHelper.dll}, rootkit driver, RAT client) decompiled separately, including identifying the helper DLL's \code{DoD3D11InstallUsingMSI} as a Rust-compiled reimplementation of \code{env::var()}. \\
\rowline
WhiteSnakeStealer & RC4 key recovered via known-plaintext analysis, and the full \code{WSR\$} C2 packet format (RC4-encrypted, Deflate-compressed system-info XML with an RSA-wrapped key) reconstructed directly from the disassembly. \\
\bottomrule
\end{tabular}
\caption{One concrete example per family of independently-recovered
evidence used to build ground truth, applying the four-step protocol
above.}
\label{tab:validation}
\end{table}

Each report's ATT\&CK Techniques table is then parsed into a structured
ground-truth file (\code{extract\_ground\_truth.py}), which the
evaluation harness compares against the pipeline's own
\code{attck\_mapping.json}.

Two matching modes are reported. \emph{Strict} matching requires an
exact technique identifier match. \emph{Family-level} matching treats
a parent technique and any of its sub-techniques as equivalent (a
pipeline finding of \technique{T1547.001} counts against a ground-truth
entry of \technique{T1547} and vice versa), on the basis that a
security analyst deciding what to do about a persistence mechanism
cares first about \emph{which} tactic and technique family is present,
with sub-technique granularity a refinement rather than a different
finding. Precision, recall, and F1 are computed under both modes for
each sample; Table~\ref{tab:results} reports family-level scores, the
stricter numbers being materially lower only where a specific
sub-technique genuinely was not distinguished.

\subsection{Results}

\begin{table}[h]
\centering
\small
\begin{tabular}{@{}lcccc@{}}
\toprule
\textbf{Sample} & \textbf{Precision} & \textbf{Recall} & \textbf{F1} & \textbf{Notes} \\
\midrule
Akira (ransomware) & 0.87 & 1.00 & 0.93 & single-stage binary \\
\rowline
WhiteSnakeStealer (.NET) & 0.97 & 0.70 & 0.81 & single-stage binary \\
\rowline
RoningLoader (parent only) & 0.95 & 0.77 & 0.85 & loader only \\
\rowline
RoningLoader (+resubmission) & 0.86 & 0.92 & 0.89 & loader + dropped components \\
\bottomrule
\end{tabular}
\caption{Family-level precision/recall/F1 per validated sample.}
\label{tab:results}
\end{table}

RoningLoader's two rows isolate exactly what the resubmission loop
(\S\ref{sec:resubmission}) contributes: scoring the loader binary alone
against ground truth written for its \emph{entire} infection chain
understates recall by construction, since the rootkit driver's and the
RAT client's own techniques are real findings the ground truth expects
but the parent binary alone cannot exhibit. Recall rises 16 points once
the three dropped components are independently analyzed and their
findings counted, at a 1-point precision cost from the added surface
area, a trade the aggregate F1 improvement (0.76~$\rightarrow$~0.84)
shows was worth taking.

Precision is high across all four rows (0.84--1.00), i.e.\ the small
number of findings the pipeline reports at all are seldom wrong. This
is the direct, intended consequence of the coherent-combination
philosophy in \S\ref{sec:static-rules} and the signature-curation
discipline in \S\ref{sec:dynamic}: both trade some recall for high
confidence that what \emph{is} reported can be trusted, which the
false-positive audit below exists to check rather than assume.

\subsection{The False-Positive Audit}
\label{sec:fp-audit}

Precision alone does not distinguish ``the pipeline is well-calibrated''
from ``the ground truth happens to be permissive.'' Unlike a reputation
aggregator that displays each engine's or sandbox's tag as given, every
discrepancy between pipeline output and ground truth found across the
three samples, 15 in total (5 Akira, 4 WhiteSnakeStealer, 6
RoningLoader), was individually traced to a specific, checkable cause
rather than accepted or dismissed on the aggregate number alone. Each
was checked against the \emph{full text} of the corresponding manual
report, not just its summary table, since a technique can be genuinely
discussed in a report's narrative without appearing in its own ATT\&CK
table.

The outcome of this audit, by cause:

\begin{itemize}
  \item \textbf{Ten genuine pipeline errors}, each fixed at its root
    cause: three static rule mis-mappings that did not match the
    claimed technique's actual MITRE definition (e.g.\
    \technique{GetSystemInfo} mapped to System Information Discovery,
    \technique{T1082}, when MITRE's own definition is specifically
    about \emph{virtualization/sandbox} detection,
    \technique{T1497}); one combination-check gap that let a single
    generic import trigger \technique{T1055} on its own; and six CAPE
    community-signature mappings whose own raw evidence field, on
    inspection, named a target inconsistent with the claimed technique
    (detailed in \S\ref{sec:dynamic}).
  \item \textbf{Two confidence-calibration corrections}, on Akira and
    RoningLoader: the same combination-check gap behind the
    \technique{T1055} fix above recurred for \technique{T1070}
    (Indicator Removal) -- \code{CreateFileW}/\code{CreateFile}, already
    commented ``generic'' in the mapping table, was still blindly
    counted as corroboration alongside \code{DeleteFileW}/\code{DeleteFile}
    and promoted to \emph{high} confidence. Migrating ground truth to
    ATT\&CK v19 (\S\ref{sec:limitations}) is what surfaced it: the old
    \technique{T1070.001} ground-truth entry had been matching this
    finding purely by ID-prefix coincidence under family-level
    matching, not because the evidence was genuinely about the entry's
    actual subject (event-log clearing); moving that entry to its
    v19 identity (\technique{T1685.005}, no longer sharing a
    \technique{T1070} prefix) broke the coincidental match. Fixed the
    same way as the \technique{T1055} case, excluding the generic
    import/string from corroboration in either direction. Unlike that
    case, the underlying evidence here is real (both binaries do
    import these functions), just not specific enough to confirm the
    claimed technique, so the finding remains present post-fix at its
    now-honest \emph{medium}/\emph{low} confidence rather than
    disappearing outright -- a direct illustration of the
    confidence/precision relationship examined in
    \S\ref{sec:confidence-calibration}, not a counterexample to it.
  \item \textbf{One ground-truth extraction gap}, not a pipeline
    error: WhiteSnakeStealer's own manual report discusses reflective
    DLL loading in its narrative text but the corresponding
    \technique{T1620} row was missing from the report's own ATT\&CK
    table, so automatic ground-truth extraction never captured it. The
    pipeline's finding was correct; the source document was fixed
    instead of the pipeline.
  \item \textbf{Two left deliberately open}, on RoningLoader
    (\technique{T1027}, \technique{T1497}): the pipeline's supporting
    evidence for both is real, but no explicit analyst confirmation
    exists anywhere in the report's text, unlike the
    \technique{T1620} case above. Editing ground truth on the
    pipeline's own evidence alone, with no independent textual
    confirmation, would be circular. Both remain open findings rather
    than resolved either way.
\end{itemize}

This asymmetry, one ground-truth gap fixed, two similar-looking
discrepancies left open, is deliberate: the distinguishing factor is
whether independent confirmation exists in the source document, not
whether the pipeline's own evidence looks convincing. The practical
payoff of running this audit at all is visible in the fixed column: 12
of 15 discrepancies turned into concrete, verifiable improvements to
the rule set (\S\ref{sec:coverage}) and the aggregation logic
(\S\ref{sec:case-installer}) rather than being absorbed silently into
an aggregate score, which is precisely the diagnostic step a
side-by-side verdict display cannot offer.

\subsection{Confidence Calibration}
\label{sec:confidence-calibration}

If the confidence reconciliation model (\S\ref{sec:reconciliation}) is
meaningful rather than cosmetic, precision should be measurably higher
among \emph{high}-confidence findings (those corroborated across
sources or across technique levels) than among single-source
\emph{medium}/\emph{low} findings. Across the three validated samples,
14 of the 15 discrepancies identified in the audit above originated
from a single-source finding. The one exception, RoningLoader's
\technique{T1497}, reached \emph{high} confidence through genuine
cross-source corroboration (a coherent anti-sandbox API combination in
the static evidence, three independent CAPE signatures in the dynamic
evidence), yet remains one of the two findings left deliberately open
above, specifically because no independent textual confirmation exists
in the manual report, not because either source's evidence is weak.
Read this way it does not contradict the calibration argument: the
model correctly recognized strong, corroborated evidence; what is
genuinely uncertain is whether the ground truth is complete, the same
asymmetry the \technique{T1620} case establishes elsewhere in this
audit, not whether the confidence tier assigned to it was earned. This
is consistent with the reconciliation model doing real calibration work
rather than merely relabeling findings, though the small sample size
(three families) means this should be read as a directional signal, not
a statistically powered claim, a limitation stated explicitly in
\S\ref{sec:limitations}.
